\section{Introduction}

Charlotte is a research prototype of an agent that is \emph{taught} how to do everyday household tasks, such as boiling an egg, fixing a flat tyre or planning a picnic. Later it uses what it was taught to help a user. The agent's memory, reasoning and replies are symbolic: knowledge is a set of Datalog facts, and every claim in an answer can be traced to a taught fact or a rule. A language model plays the teacher and, where the agent cannot yet read a sentence, a fallback reader. It never writes the agent's answers.

The project asks three questions:
\begin{enumerate}
  \item Is knowledge that is taught once reused correctly later, when the user words a request differently?
  \item Does the agent transfer what it learned to new situations (a missing tool, an extra ingredient), to new tasks and to a whole new domain?
  \item How much of the work of reading English can the agent learn, instead of having it written into its code?
\end{enumerate}

This paper reports three iterations. Version 1 built a symbolic core with a hand-written parser and found that it could not represent procedures at all. Version 2.0 added procedural knowledge, an LLM translator and a dialogue curriculum. It became fluent, but its gains came from about 225 English strings hard-coded into the agent and from polishing replies for style judges. Design 3.0 changes what is measured and where knowledge lives:
\begin{itemize}
  \item answers are lists of facts and relations, checked exactly against a re-authored ground truth;
  \item the agent's language layer is a learned model with the teacher as its fallback;
  \item adaptation to new situations is learned from cases, and gated by evidence;
  \item a sealed held-out set of \hThreeS{} new-situation, \hThreeT{} new-task and \hThreeG{} new-domain scenarios measures the result.
\end{itemize}
Every number in this paper is generated from the repository's data by a script.
